\documentclass[11pt, a4paper]{article}

\usepackage{amsmath, amssymb, amsthm}
\usepackage{bm}
\usepackage{geometry}
\usepackage{booktabs}
\usepackage{hyperref}
\usepackage{xcolor}
\usepackage{microtype}
\usepackage{parskip}

\geometry{margin=2.5cm}

\hypersetup{
  colorlinks = true,
  linkcolor  = blue!60!black,
  urlcolor   = blue!60!black,
}

% Convenience macros
\newcommand{\x}{\bm{x}}
\newcommand{\z}{\bm{z}}
\newcommand{\y}{\bm{y}}
\newcommand{\w}{\bm{w}}
\newcommand{\vv}{\bm{v}}
\newcommand{\E}{\mathbb{E}}
\newcommand{\N}{\mathcal{N}}
\newcommand{\R}{\mathbb{R}}
\newcommand{\given}{\mid}

\theoremstyle{remark}
\newtheorem*{remark}{Remark}

% ─────────────────────────────────────────────────────────────────────────────

\title{\textbf{GPS Track Smoothing via the Kalman Filter\\and RTS Backward Smoother}\\[0.4em]
       \large A Statistical Derivation}
\author{}
\date{}

\begin{document}
\maketitle
\tableofcontents
\bigskip

% ─────────────────────────────────────────────────────────────────────────────
\section{Overview}

Raw GPS recordings of a nordic ski track contain two types of unwanted
signal: high-frequency \emph{jitter} (Gaussian-like measurement noise of
a few metres) and occasional hard \emph{spikes} (outlying fixes tens or
hundreds of metres away from the true path).  A plain smoothing spline
treats every observation symmetrically and is therefore pulled toward
spikes.  The approach taken here instead models the data-generating
process explicitly as a \emph{linear Gaussian state-space model} and
computes the exact posterior over the latent track using the
\textbf{Kalman filter} (forward pass) and the \textbf{Rauch--Tung--Striebel
(RTS) smoother} (backward pass), with a \emph{Mahalanobis innovation gate}
to reject hard outliers before they enter the filter.

% ─────────────────────────────────────────────────────────────────────────────
\section{State-Space Model}

\subsection{Coordinate system}

Each GPS segment is worked in a local flat-Earth Cartesian frame centred on
the first point of that segment.  Latitude and longitude are converted to
metres via
\begin{align}
  x &= (\lambda - \lambda_0)\cdot 111{,}319.9\cos\phi_0, \\
  y &= (\phi   - \phi_0)   \cdot 111{,}319.9,
\end{align}
where $\phi_0,\lambda_0$ are the reference latitude and longitude in
radians and degrees, respectively.  This approximation is accurate to
better than 0.1\,\% for segment lengths below 50\,km.

\subsection{Latent state and observations}

At each time step $t$ define the four-dimensional latent state
\begin{equation}
  \x_t = \begin{pmatrix} x_t \\ y_t \\ \dot x_t \\ \dot y_t \end{pmatrix}
  \in \R^4,
\end{equation}
collecting horizontal position $(x_t, y_t)$ and velocity
$(\dot x_t, \dot y_t)$ in metres and metres per second, respectively.
The GPS receiver produces the two-dimensional observation
\begin{equation}
  \z_t = \begin{pmatrix} x_t^{\mathrm{GPS}} \\ y_t^{\mathrm{GPS}} \end{pmatrix}
  \in \R^2.
\end{equation}

\subsection{Transition equation}

We adopt a \emph{constant-velocity model} perturbed by random accelerations:
\begin{equation}
  \x_t = F_t\,\x_{t-1} + \w_t, \qquad \w_t \sim \N(\bm{0},\, Q_t),
  \label{eq:transition}
\end{equation}
where $\Delta t_t = t - (t-1)$ is the elapsed time between consecutive
GPS fixes (in seconds) and
\begin{equation}
  F_t = \begin{pmatrix}
    1 & 0 & \Delta t_t & 0 \\
    0 & 1 & 0 & \Delta t_t \\
    0 & 0 & 1 & 0 \\
    0 & 0 & 0 & 1
  \end{pmatrix}.
\end{equation}

The process noise $\w_t$ arises from unmodelled accelerations.
Treating acceleration as white noise with standard deviation
$\sigma_a$ and integrating twice yields the \emph{continuous Wiener
process acceleration} (CWPA) noise covariance:
\begin{equation}
  Q_t = \sigma_a^2
  \begin{pmatrix}
    \Delta t_t^4/4 & 0              & \Delta t_t^3/2 & 0 \\
    0              & \Delta t_t^4/4 & 0              & \Delta t_t^3/2 \\
    \Delta t_t^3/2 & 0              & \Delta t_t^2   & 0 \\
    0              & \Delta t_t^3/2 & 0              & \Delta t_t^2
  \end{pmatrix}.
  \label{eq:Q}
\end{equation}

\subsection{Measurement equation}

Only position is observed:
\begin{equation}
  \z_t = H\,\x_t + \vv_t, \qquad \vv_t \sim \N(\bm{0},\, R),
  \label{eq:measurement}
\end{equation}
where $H = [I_2 \mid \bm{0}_{2\times 2}]$ selects the position components
and $R = \sigma_{\mathrm{GPS}}^2\,I_2$ is an isotropic measurement noise
covariance.  Both noise sequences $\{\w_t\}$ and $\{\vv_t\}$ are mutually
independent and independent over time, and are independent of the initial
state.

Under these assumptions, \eqref{eq:transition}--\eqref{eq:measurement}
define a \textbf{linear Gaussian state-space model}.  The joint
distribution over $(\x_{0:T}, \z_{1:T})$ is multivariate Gaussian, and
all conditional distributions of interest are Gaussian with closed-form
parameters.

% ─────────────────────────────────────────────────────────────────────────────
\section{Forward Kalman Filter}

The Kalman filter computes the \emph{filtering distribution}
$p(\x_t \given \z_{1:t}) = \N(\hat\x_{t|t},\, P_{t|t})$ recursively.

\subsection{Initialisation}

We initialise at $t=0$ (the first point of the segment) with position
taken from the GPS fix and velocity set to zero:
\begin{equation}
  \hat\x_{0|0} = \begin{pmatrix} z_{0,x} \\ z_{0,y} \\ 0 \\ 0 \end{pmatrix},
  \qquad
  P_{0|0} = \operatorname{diag}\!\left(\sigma_{\mathrm{GPS}}^2,\,
    \sigma_{\mathrm{GPS}}^2,\, \sigma_{v,0}^2,\, \sigma_{v,0}^2\right),
\end{equation}
where $\sigma_{v,0} = 3$\,m/s is a vague prior on initial speed.

\subsection{Predict step}

Marginalising over $\x_{t-1}$ using the transition \eqref{eq:transition}
gives the \emph{one-step-ahead predictive distribution}
$p(\x_t \given \z_{1:t-1}) = \N(\hat\x_{t|t-1},\, P_{t|t-1})$:
\begin{align}
  \hat\x_{t|t-1} &= F_t\,\hat\x_{t-1|t-1}, \label{eq:predict_mean}\\
  P_{t|t-1}      &= F_t\,P_{t-1|t-1}\,F_t^\top + Q_t. \label{eq:predict_cov}
\end{align}

\subsection{Innovation}

Define the \emph{innovation} (the prediction error in observation space):
\begin{equation}
  \y_t = \z_t - H\,\hat\x_{t|t-1},
\end{equation}
with innovation covariance
\begin{equation}
  S_t = H\,P_{t|t-1}\,H^\top + R.
\end{equation}
Under the model, $\y_t \sim \N(\bm{0}, S_t)$ when $\z_t$ is consistent
with the predicted state.

\subsection{Update step}

Conditioning $p(\x_t \given \z_{1:t-1})$ on $\z_t$ via Bayes' rule gives
\begin{align}
  \hat\x_{t|t} &= \hat\x_{t|t-1} + K_t\,\y_t, \label{eq:update_mean}\\
  P_{t|t}      &= (I - K_t H)\,P_{t|t-1},       \label{eq:update_cov}
\end{align}
where the \emph{Kalman gain}
\begin{equation}
  K_t = P_{t|t-1}\,H^\top\,S_t^{-1}
\end{equation}
is the matrix of regression coefficients of $\x_t$ on $\y_t$ under the
joint predictive distribution.  Equivalently, $K_t$ minimises
$\operatorname{tr}(P_{t|t})$, making the posterior variance as small
as possible.

\begin{remark}
The update \eqref{eq:update_mean}--\eqref{eq:update_cov} is the standard
formula for conditioning a joint Gaussian.  If we write
$\binom{\x_t}{\z_t} \sim \N(\cdot, \Sigma)$ under the prior, the
posterior mean is the prior mean plus the cross-covariance times the
precision of $\z_t$ times the residual, which is exactly $K_t \y_t$.
\end{remark}

% ─────────────────────────────────────────────────────────────────────────────
\section{Innovation Gating}

Before the update step, each incoming GPS fix is subjected to an outlier
test.  Under the model, the squared Mahalanobis distance of the innovation
\begin{equation}
  d_t^2 = \y_t^\top\,S_t^{-1}\,\y_t
  \label{eq:mahal}
\end{equation}
follows a chi-squared distribution with $m = 2$ degrees of freedom
(the dimension of $\z_t$) when the observation is genuine:
\begin{equation}
  d_t^2 \;\big|\; \text{inlier} \;\sim\; \chi^2_2.
\end{equation}

A fix is \emph{gated out} if $d_t^2 > g$, where $g$ is a threshold chosen
from the $\chi^2_2$ quantile function.  Common choices are given in
Table~\ref{tab:gate}.  When a fix is rejected the update step is skipped:
$\hat\x_{t|t} = \hat\x_{t|t-1}$ and $P_{t|t} = P_{t|t-1}$.  The filter
coasts on the dynamics model until the next accepted fix.

\begin{table}[h]
\centering
\caption{Chi-squared gate thresholds for $\chi^2_2$.}
\label{tab:gate}
\begin{tabular}{lcc}
\toprule
Coverage & $g = \chi^2_2(p)$ & Effective $\sigma$-radius \\
\midrule
95.0\,\% &  5.99 & ${\approx}2.45\,\sigma_{\mathrm{eff}}$ \\
99.0\,\% &  9.21 & ${\approx}3.03\,\sigma_{\mathrm{eff}}$ \\
99.9\,\% & 13.82 & ${\approx}3.72\,\sigma_{\mathrm{eff}}$ \\
\bottomrule
\end{tabular}
\end{table}

The default threshold $g = 13.82$ (99.9\,\%) is deliberately conservative:
only measurements that are extraordinarily far from the predicted position
are rejected outright, because a skier can in principle accelerate or
change direction rapidly.  The smoother handles mild GPS noise without
needing to gate it.

% ─────────────────────────────────────────────────────────────────────────────
\section{RTS Backward Smoother}

The forward Kalman filter is a \emph{causal} filter: $\hat\x_{t|t}$
conditions only on $\z_{1:t}$.  The RTS smoother computes the full
posterior
\begin{equation}
  p(\x_t \given \z_{1:T}) = \N(\hat\x_{t|T},\, P_{t|T})
\end{equation}
for all $t$ simultaneously, by making a single backward pass through the
stored forward-filter output.

\subsection{Derivation sketch}

By the Markov property of the state-space model,
\begin{equation}
  p(\x_t \given \z_{1:T})
  = \int p(\x_t \given \x_{t+1},\, \z_{1:t})\;
         p(\x_{t+1} \given \z_{1:T})\; \mathrm{d}\x_{t+1}.
\end{equation}
The integrand is Gaussian in both factors.  The joint distribution of
$(\x_t, \x_{t+1})$ given $\z_{1:t}$ is
\begin{equation}
  \begin{pmatrix} \x_t \\ \x_{t+1} \end{pmatrix} \Bigg|\, \z_{1:t}
  \;\sim\;
  \N\!\left(
    \begin{pmatrix} \hat\x_{t|t} \\ \hat\x_{t+1|t} \end{pmatrix},\;
    \begin{pmatrix} P_{t|t} & P_{t|t} F_{t+1}^\top \\
                    F_{t+1} P_{t|t} & P_{t+1|t} \end{pmatrix}
  \right).
\end{equation}
Conditioning $\x_t$ on $\x_{t+1}$ (standard Gaussian conditioning formula)
and then marginalising over $\x_{t+1}$ using the already-computed smoothing
distribution at $t+1$ yields the \textbf{RTS recursion}:
\begin{align}
  G_t &= P_{t|t}\,F_{t+1}^\top\,P_{t+1|t}^{-1}, \label{eq:gain_rts}\\[4pt]
  \hat\x_{t|T} &= \hat\x_{t|t}
                  + G_t\!\left(\hat\x_{t+1|T} - \hat\x_{t+1|t}\right),
                  \label{eq:smooth_mean}\\[4pt]
  P_{t|T}      &= P_{t|t}
                  + G_t\!\left(P_{t+1|T} - P_{t+1|t}\right)G_t^\top,
                  \label{eq:smooth_cov}
\end{align}
initialised with $\hat\x_{T|T}$ and $P_{T|T}$ from the forward filter
(the final filtering distribution already equals the smoothing distribution
at time $T$).

\subsection{Interpretation}

$G_t$ in \eqref{eq:gain_rts} is the cross-covariance of $\x_t$ and
$\x_{t+1}$ (under the forward predictive distribution) scaled by the
precision of $\x_{t+1}$.  It is the \emph{backward analogue} of the
Kalman gain.

Equation~\eqref{eq:smooth_mean} says: take the forward-filtered estimate
$\hat\x_{t|t}$, compute how much the smoother revised $\hat\x_{t+1}$
relative to the filter's prediction ($\hat\x_{t+1|T} - \hat\x_{t+1|t}$),
and propagate that revision backwards through $G_t$.  If a clean GPS fix
at time $t+3$ showed the skier 50\,m further north than predicted, the
smoother pulls estimates at $t$, $t+1$, and $t+2$ northward as well,
weighted by the smoother gain at each step.

This is why the RTS smoother is fundamentally different from a one-sided
exponential moving average or a Savitzky--Golay filter: future observations
actively improve past estimates.

% ─────────────────────────────────────────────────────────────────────────────
\section{Hyperparameters}

The model has three hyperparameters.  All are interpretable in physical
units and can be tuned from domain knowledge or, in principle, by maximum
marginal likelihood (EM algorithm on the innovations).

\subsection{\texorpdfstring{$\sigma_{\mathrm{GPS}}$}{sigma\_GPS} — GPS measurement noise}

\textbf{Default:} 5\,m.

$\sigma_{\mathrm{GPS}}$ is the assumed one-sigma horizontal accuracy of the
GPS receiver.  It enters the measurement covariance $R = \sigma_{\mathrm{GPS}}^2 I_2$
and therefore the innovation covariance $S_t$ and Kalman gain $K_t$.

\begin{itemize}
  \item \textbf{Smaller $\sigma_{\mathrm{GPS}}$}: the filter trusts GPS fixes more.
        The smoothed track stays closer to the raw observations.
        Appropriate when the receiver is high-quality and conditions are
        good (open terrain, strong satellite geometry).

  \item \textbf{Larger $\sigma_{\mathrm{GPS}}$}: the filter trusts GPS fixes less
        and relies more on the dynamics model.
        The smoothed track is more regular but may cut corners.
        Appropriate in forests or deep valleys where multipath is common.
\end{itemize}

The Kalman gain satisfies $K_t \to I$ as $\sigma_{\mathrm{GPS}} \to 0$
(fully trust GPS) and $K_t \to 0$ as $\sigma_{\mathrm{GPS}} \to \infty$
(ignore GPS, coast on the model).

\subsection{\texorpdfstring{$\sigma_a$}{sigma\_a} — process noise / expected acceleration}

\textbf{Default:} 1.5\,m/s$^2$.

$\sigma_a$ characterises how much the skier's velocity is expected to
change between consecutive fixes, and is the standard deviation of the
latent acceleration.  It scales the entire process noise matrix $Q_t$
in \eqref{eq:Q}.

\begin{itemize}
  \item \textbf{Smaller $\sigma_a$}: the model expects nearly constant velocity.
        Turns and acceleration are penalised — the track is very smooth
        but tight bends may be cut short.
        Analogous to a large smoothing penalty $\lambda$ in a spline.

  \item \textbf{Larger $\sigma_a$}: the model allows rapid changes in speed and
        direction.  The smoother is more willing to follow the GPS closely.
        Analogous to a small $\lambda$.
\end{itemize}

The trade-off between smoothing and fidelity is governed by the
\emph{ratio} $\sigma_a / \sigma_{\mathrm{GPS}}$.  A useful rule of thumb:
for a 1-second logging interval, $\sigma_a = 1$--$2$\,m/s$^2$ covers
realistic skiing accelerations (typical values are 0.3--1.5\,m/s$^2$
on groomed trails).  If the logging rate is lower (e.g.\ every 5\,s),
$Q_t$ grows with $\Delta t^2$ so the model naturally allows more freedom
between fixes without changing $\sigma_a$.

\subsection{\texorpdfstring{$g$}{g} — innovation gate threshold}

\textbf{Default:} 13.82\;($= \chi^2_2(0.999)$).

$g$ is the chi-squared threshold used in the Mahalanobis gate test
\eqref{eq:mahal}.  A fix is \emph{rejected} if $d_t^2 > g$.

\begin{itemize}
  \item \textbf{Smaller $g$} (e.g.\ 5.99 at 95\,\%): the gate is tighter.
        More fixes are rejected and the smoother coasts on the model more
        often.  Useful when the GPS is known to produce frequent hard
        spikes, but risks discarding genuine rapid movements.

  \item \textbf{Larger $g$} (e.g.\ 20): almost no fixes are rejected.
        The smoother itself handles mild outliers by down-weighting them
        via $K_t$ — but it cannot fully suppress a measurement 200\,m off
        track.  The gate provides a hard safety net.
\end{itemize}

The default 99.9\,\% gate is deliberately permissive because the filter's
own weighting already attenuates mild GPS noise.  Only measurements that
are extremely improbable under the current motion model are discarded.

\subsection{Summary table}

\begin{table}[h]
\centering
\caption{Hyperparameter reference.}
\begin{tabular}{llll}
\toprule
Parameter & Default & Unit & Effect of increasing \\
\midrule
$\sigma_{\mathrm{GPS}}$ & 5.0   & m        & Less trust in GPS; smoother track \\
$\sigma_a$              & 1.5   & m/s$^2$  & More dynamics allowed; track follows GPS more closely \\
$g$                     & 13.82 & ---      & Fewer points rejected; more smoothed by the filter \\
\bottomrule
\end{tabular}
\end{table}

% ─────────────────────────────────────────────────────────────────────────────
\section{Comparison with Smoothing Splines}

A cubic smoothing spline minimises
\begin{equation}
  \sum_{t=1}^T \bigl(f(t_t) - z_t\bigr)^2
  + \lambda \int \bigl[f''(u)\bigr]^2\, \mathrm{d}u,
  \label{eq:spline}
\end{equation}
penalising curvature globally.  The differences from the Kalman smoother
are:

\begin{enumerate}
  \item \textbf{Motion model.}
        The spline has no concept of velocity or physically plausible
        acceleration.  The Kalman model encodes Newton's laws: position
        changes at the rate of velocity, and velocity changes due to
        acceleration.  This matters at the ends of the track and across
        gaps caused by gated outliers.

  \item \textbf{Outlier robustness.}
        The spline objective \eqref{eq:spline} is a sum of squared
        residuals; a single 200\,m spike distorts the fit over a wide
        neighbourhood.  The Kalman gate removes such measurements entirely
        before the smoother sees them.

  \item \textbf{Variable time steps.}
        The spline in \eqref{eq:spline} treats observations as equally
        spaced unless the knot structure is adapted.  The Kalman model
        uses the actual timestamp differences $\Delta t_t$ in both $F_t$
        and $Q_t$, so it naturally handles irregular sampling.

  \item \textbf{Two-dimensional consistency.}
        Applying a spline independently to latitude and longitude can
        produce paths that violate the joint dynamics (e.g.\ an
        impossibly sharp change in the north--south direction while
        east--west is smooth).  The Kalman model couples $x$ and $y$
        through the shared velocity state and covariance matrix.

  \item \textbf{Uncertainty quantification.}
        The Kalman smoother produces $P_{t|T}$, a full covariance matrix
        for each smoothed point.  This can be used to construct credible
        intervals on the position or speed at any time step.
        The smoothing spline provides no such uncertainty.
\end{enumerate}

\end{document}
